% ===== 底本 第 3 章《环簇 Fano 集的 SYZ 镜像对称》；片段 tr03b；源块 B116--B131 =====

% [源 B116]
\subsection{同调镜像对称}

% [源 B117]
在深入 SYZ 镜像构造之前，我们先简要说明环簇的同调镜像对称（HMS）的已知结果。最简单的环簇例子是带有线性 $\mathbb{T}^n$-作用的代数环面 $(\mathbb{C}^{\ast})^n$。它是阿贝尔李代数 $\mathbb{R}^n=\mathfrak{t}^n$ 上的正则拉格朗日纤维化。

% [源 B118]
\begin{theorem*}
$\left(\mathbb{C}^{\ast}\right)^n$ 是\emph{自镜像的}（self-mirror）。换言之，存在范畴等价
\[
  D^b\operatorname{Coh} \left(\left(\mathbb{C}^{\ast}\right)^n\right)\cong D^{\pi}\mathcal{W} \left(\left(\mathbb{C}^{\ast}\right)^n\right).
\]
\end{theorem*}

% [源 B119]
\begin{proof}
由 [Abo11] 我们知道，包裹 Fukaya 范畴 $\mathcal{W}\left(\mathbb{C}^{\ast}\right)^n$ 可视为 $\mathbb{T}^n$ 的余切丛，它由单个余切纤维 $T^{\ast}_q\mathbb{T}^n$ 生成，其中 $q\in\mathbb{T}^n$。于是剩下的只需证明存在 $A_{\infty}$-代数同构
\[
  \operatorname{End}_{D^{\pi}\mathcal{W} \left(\left(\mathbb{C}^{\ast}\right)^n\right)} \left(T^{\ast}_q\mathbb{T}^n\right)\cong \operatorname{End}_{D^b\operatorname{Coh} \left(\left(\mathbb{C}^{\ast}\right)^n\right)} \left(\mathcal{O}_{\left(\mathbb{C}^{\ast}\right)^n}\right).
\]
由 [Abo12] 我们知道，左端由环路空间 $\Omega_q\mathbb{T}^n$ 的同调给出，而同构于群代数 $\mathbb{C}[\pi_1 (\mathbb{T}^n)]\cong\mathbb{C}[x_1^{\pm 1},\dotsc, x_n^{\pm 1}]\cong\Gamma\left(\left(\mathbb{C}^{\ast}\right)^n,\mathcal{O}\right)$。
\end{proof}

% [Tutor Note: B119 末句的 $\mathbb{C}[x_1^{\pm 1},\dotsc, x_n^{\pm 1}]$ 与 $\Gamma\left(\left(\mathbb{C}^{\ast}\right)^n,\mathcal{O}\right)$ 是同一个环：环面上正则函数恰为 Laurent 多项式。可用 $n=1$ 核验：$(\mathbb{C}^{\ast})^1$ 上正则函数正是 $\mathbb{C}[x^{\pm 1}]$。]

% [源 B120]
这个定理还告诉我们，$T^{\ast}\mathbb{T}^n$ 的余切纤维镜像于 $\left(\mathbb{C}^{\ast}\right)^n$ 的结构层（structure sheaf），因此我们可以尝试从这个同调镜像等价中恢复出所有的点对象，这就导向 SYZ 构造。注意由于 $\mathbb{T}^n$ 是李群，$T^{\ast}\mathbb{T}^n\cong\mathbb{T}^n\times\mathbb{R}^n$ 是平凡的，于是由乘积分解我们得到一族拉格朗日环面 $\{L_b\}_{b\in \mathbb{R}^n}$。

% [Tutor Note: B120 原文 "decompostion" 拼写有误，应为 decomposition（分解）；译文按「分解」处理，数学内容不变，已写入回执。]

% [源 B121]
下面考察这些拉格朗日环面的拉格朗日 Floer 理论。首先容易看出，当 $b\neq b'$ 时 $HF^{\ast} (L_b, L_{b'})=0$，因为它们没有几何交点；而 $HF^{\ast} (L_b,T^{\ast}_q\mathbb{T}^n)\cong\mathbb{R}$，因为它们有唯一的横截交点。唯一有趣的复形是 $CF^{\ast} (L_b,L_b)$（$b\in\mathbb{R}^n$）。

% [源 B122]
回忆 $T^{\ast}\mathbb{T}^n$ 上的标准辛结构由 $-d\lambda_{can}$ 给出，其中 $\lambda_{can}=\sum_{i=1}^n (\log r_i)d\theta_i$，这里 $(\log r_i)_{i=1}^n$ 是余切坐标，$(\theta_i)$ 是环面坐标。回忆如下定义：

% [源 B123]
\begin{definition*}
恰当辛流形 $(M,\omega=d\lambda)$ 的拉格朗日子流形 $L\subseteq (M,\lambda)$ 称为\emph{恰当的}（exact），如果 $\lambda\vert_L$ 是 $L$ 上的恰当 $1$-形式。
\end{definition*}

% [源 B124]
拉格朗日子流形 $L$ 的恰当性阻止了以 $L$ 为边界的任何非常值伪全纯圆盘的存在，因此 Floer 上同调容易计算：它与 Morse 上同调相同，于是对恰当的拉格朗日环面 $L_b$ 有
\[
  HF^{\ast} (L_b,L_b)\cong H^{\ast} (L_b;\mathbb{R})
\]

% [源 B125]
然而不难验证，$T^{\ast}\mathbb{T}^n$ 中唯一的恰当拉格朗日环面 $L_b$ 是零截面 $L_0$。对其余所有拉格朗日环面 $L_b$（$b\neq 0$），以 $L_b$ 为边界的全纯圆盘的面积，由 $L_b$ 中边界曲线类与 Liouville 形式 $\lambda$ 的积分决定；更具体地说，由该曲线类与 $\lambda$ 限制到环面 $L_b$ 上所得的 $1$-形式 $\lambda\vert_{L_b}$ 的积分决定。

% [源 B126]
我们把 $A_{\infty}$-结构写出来，以看出这些非平凡全纯曲线的作用：记 $CF^{\ast} (L_b,L_b)=\Omega(L_b)$ 为 $L_b$ 的 de Rham 代数，则 $A_{\infty}$-结构 $\mathfrak{m}^d\colon CF^{\ast} (L_b,L_b)\otimes\dotsb\otimes CF^{\ast} (L_b, L_b)\to CF^{\ast} (L_b,L_b)[2-d]$ 由
\[
  \resizebox{0.985\textwidth}{!}{$\displaystyle \mathfrak{m}^d (x_1,\dotsb, x_d)=\mu^d (x_1,\dotsb, x_d)+\sum_{\beta}\sum_{\deg (x_0)=\sum\deg (x_i)+2-k+\mu(\beta)}\left(\int_{[\mathcal{M}_{d+1} (\beta)]^{vir}}\ev_0^{\ast} x_0\cup\ev_1^{\ast} x_1\cup\dotsb\cup\ev_d^{\ast} x_d\right)T^{\lambda (\partial\beta)} x_0, $}
\]
给出，其中 $\mu^d$ 是 de Rham 代数 $\Omega^{\ast} (L_0)$ 上的 Morse $A_{\infty}$-结构，$\beta\in\pi_2 (X,L_b)$ 是圆盘类，$\partial\beta\in\pi_1 (L_b)$ 是相应的边界类。注意值 $\lambda(\partial\beta)$ 只依赖于同调类 $\left[\lambda\vert_{L_b}\right]\in H^1 (L_b;\mathbb{R})\cong H^1 (L_0;\mathbb{R})\cong\mathbb{R}^2$ 与 $[\partial\beta]\in\pi_1 (L_b)\cong\pi_1 (L_0)\cong\mathbb{Z}^2$，于是我们看到，$A_{\infty}$-代数 $CF^{\ast} (L_b,L_b)$ 是 $A_{\infty}$-代数 $CF^{\ast} (L_0,L_0)$ 沿元素 $[\lambda_b]\in H^1_{dR} (L_0;\mathbb{R})$ 的形变，其中 $CF^{\ast}$ 的生成元相同，而 $A_{\infty}$-结构形变为 $\mathfrak{m}^d_b$。

% [Tutor Note: B126 的大公式里内层求和条件是 $\deg (x_0)=\sum\deg (x_i)+2-k+\mu(\beta)$，但全段只出现 $\mathfrak{m}^d$ 与 $\mu^d$，从未引入 $k$；按 $A_{\infty}$-结构 $\mathfrak{m}^d$ 的次数为 $2-d$，此处应为 $\deg (x_0)=\sum\deg (x_i)+2-d+\mu(\beta)$，即 $k$ 疑为 $d$ 之笔误。照原文排印，已写入回执。]

% [Tutor Note: B126 中 $H^1 (L_b;\mathbb{R})\cong H^1 (L_0;\mathbb{R})\cong\mathbb{R}^2$ 与 $\pi_1 (L_b)\cong\pi_1 (L_0)\cong\mathbb{Z}^2$ 对一般 $n$ 而言应为 $\mathbb{R}^n$ 与 $\mathbb{Z}^n$（$L_0=\mathbb{T}^n$，$H^1(\mathbb{T}^n;\mathbb{R})\cong\mathbb{R}^n$、$\pi_1(\mathbb{T}^n)\cong\mathbb{Z}^n$）；原文写 $2$ 疑为沿用了 $n=2$ 的记法，照原文排印，已写入回执。]

% [源 B127]
在此情形，从拓扑上不难看出该形变实际上是平凡的：由于 $T^{\ast}\mathbb{T}^n\to\mathbb{R}^n$ 是纤维化，可得 $\pi_2 (T^{\ast}\mathbb{T}^n,L_b)\cong\ker (\pi_1 (L_b)\to\pi_1 (T^{\ast}\mathbb{T}^n))=0$，因此以 $L_b$ 为边界的伪全纯圆盘必为平凡圆盘。于是所有代数 $CF^{\ast} (L_b,L_b;\mathbb{R})$ 都可以典范地等同于其 de Rham 代数，带有平凡的 Morse $A_{\infty}$-结构。特别地，它们都对应镜像 $\left(\mathbb{C}^{\ast}\right)^n$ 中的光滑闭点；并且我们看到，所有配对 $(L_b,\nabla_b)$ 构成的模空间同构于 $\left(\mathbb{C}^{\ast}\right)^n$，其中 $\nabla_b$ 是 $L_b$ 上平凡复线束的平坦联络。

% [Tutor Note: B127 原文 "falt connection" 拼写有误，应为 flat connection（平坦联络）；译文按「平坦联络」处理，数学内容不变，已写入回执。]

% [源 B128]
现在转向一般的射影环簇。由于所有环簇都包含一个 Zariski 稠密开代数环面，因此所有环簇都必须是 $\left(\mathbb{C}^{\ast}\right)^n$ 的紧化。在这些情形下，我们仍得到一族以 $\left(\mathbb{C}^{\ast}\right)^n$ 为参数、带有局部系统的拉格朗日环面；但若现在考察 $CF^{\ast} (L_b,L_b)$ 的 $A_{\infty}$-结构，就可能存在以 $L_b$ 为边界、与环簇边界除子相交的伪全纯圆盘。正如从公式所见，额外圆盘的存在对应先前 $A_{\infty}$-结构的形变，因此我们预期环簇的 Fukaya 范畴是 $\left(\mathbb{C}^{\ast}\right)^n$ 的 Fukaya 范畴的形变。镜像的代数范畴也应有同样的行为，我们将在下一节回到这一点。

% [源 B129]
如果把辛结构与代数结构的角色互换，那么一般环簇栈（toric stacks）的同调镜像对称已由 Fang--Liu--Treumann--Zaslow [FLTZ11] 与 Kuwagaki [Kuw20] 用\emph{可构造层}（constructible sheaves）证明。镜像一侧的相应陈述通常称为\emph{凝聚-可构造对应}（coherent-constructible correspondence）。

% [源 B130]
\begin{theorem*}[Fang--Liu--Treumann--Zaslow '11, Kuwagaki '20]
设 $X_{\Sigma}$ 是由扇 $\Sigma\subseteq N_{\mathbb{R}}$ 相伴的环簇，则存在范畴等价
\[
  D^b\operatorname{Coh} (X_{\Sigma})\cong D^b\operatorname{Sh}_c (\mathbb{T}^n,\Lambda_{\Sigma})
\]
其中 $\Lambda_{\Sigma}\subseteq T^{\ast} T^n$ 是由扇 $\Sigma$ 决定的锥状拉格朗日子流形，而 $\operatorname{Sh}_c (\mathbb{T}^n, \Lambda_{\Sigma})$ 是奇异支撑落在 $\Lambda_{\Sigma}$ 上的可构造层构成的范畴。
\end{theorem*}

% [Tutor Note: B130 的 $\Lambda_{\Sigma}\subseteq T^{\ast} T^n$ 中写作 $T^n$，而同段及上下文其余各处写作 $\mathbb{T}^n$，两者是同一个环面；底本记号不统一，照原文保留，勿据上下文改写。]

% [源 B131]
当 $X_{\Sigma}$ 是 Fano 环簇时，[GPS24] 证明了 $\operatorname{Sh}_c (\mathbb{T}^n,\Lambda_{\Sigma})$ 等价于部分包裹 Fukaya 范畴 $\mathcal{W}\left(\left(\mathbb{C}^{\ast}\right)^n,W_{\Sigma}\right)$，其中 $W_{\Sigma}$ 是所谓的\emph{Hori-Vafa 超势}，它也由 $\Sigma$ 的组合结构决定。
